\subsection*{3.2 Trading Space for Time}
In this section, we will modify Algorithm1 to make $m$ dramatically smaller, in exchange for making the branching program longer.
Before jumping into the proof of Theorem11, we will address the main innovation of our work, namely trading off time for space. Namely we begin by building a register program that takes time $n2^{n+1}$ but uses only the $n+1$ registers $R_1^{\mathrm{in}}\ldots R_n^{\mathrm{in}},R_{\mathrm{out}}$. This is similar to what Theorem11 guarantees when $k=n$.
As in Sections3.1 and3.3, let $p_f\in\mathbb F_2[x_1,\ldots,x_n]$ be $f$ as a polynomial, let $q_f\in\mathbb F_2[\tau_1^{\mathrm{in}},\ldots,\tau_n^{\mathrm{in}},y_1,\ldots,y_n]$ be equal to $p_f$ so long as $y_i=\tau_i^{\mathrm{in}}+x_i$ for all $i$, and let $c_{S,S'}$ be the coefficient of $(\prod_{i\in S}\tau_i^{\mathrm{in}})(\prod_{i\in S'}y_i)$ in $q_f$. We define a small generalization of ToggleInput, where we can choose to toggle only a subset of our inputs:
\paragraph{ToggleSomeInputs($S'$).}
\begin{enumerate}[label=\arabic*:]
\item \textbf{procedure} ToggleSomeInputs($S'$)
\item \textbf{for} $i\in S'$ \textbf{do}
\item $R_i^{\mathrm{in}}\leftarrow R_i^{\mathrm{in}}+x_i$
\item \textbf{end for}
\item \textbf{end procedure}
\end{enumerate}
Using ToggleSomeInputs($S'$), we can replace the register $R_{S'}$ in Algorithm1 with a separate set of instructions that computes the corresponding term of $q_f$:
\paragraph{Algorithm2: A slow algorithm for computing $q_f$.}
\begin{enumerate}[label=\arabic*:]
\item \textbf{for} $S'\subseteq[n]$ \textbf{do}
\item ToggleSomeInputs($S'$)
\item $R_{\mathrm{out}}\leftarrow R_{\mathrm{out}}+\sum_{S\subseteq[n]}c_{S,S'}\cdot\prod_{i\in S\cup S'}R_i^{\mathrm{in}}$
\item ToggleSomeInputs($S'$)
\item \textbf{end for}
\end{enumerate}
Whenever Line3 is executed, $R_i^{\mathrm{in}}=y_i$ for $i\in S'$, and $R_i^{\mathrm{in}}=\tau_i^{\mathrm{in}}$ for $i\notin S'$. By construction of $q_f$, $S$ and $S'$ are disjoint whenever $c_{S,S'}\neq0$, from which it follows that $R_i^{\mathrm{in}}=\tau_i^{\mathrm{in}}$ for $i\in S$. Thus the effect of Line3 is to add $\sum_{S\subseteq[n]}c_{S,S'}(\prod_{i\in S}\tau_i^{\mathrm{in}})(\prod_{i\in S'}y_i)$ to $R_{\mathrm{out}}$. Since this is run for every subset $S'$, the overall effect of the program is to add
\[\sum_{S,S'\subseteq[n]}c_{S,S'}\left(\prod_{i\in S}\tau_i^{\mathrm{in}}\right)\left(\prod_{i\in S'}y_i\right)=p_f(x_1\ldots x_n)\]
to $R_{\mathrm{out}}$.
We now give an overview of the full proof. Our goal is to balance Algorithms1 and2 by removing the registers $R_S$ corresponding to large subsets $S$ and using slow multiplication to build the polynomial $q_f$ from the remaining small subsets. In particular, if we divide the input bits into $k$ groups each of size $\lceil n/k\rceil$, and only store all subsets within each group, then any monomial $c_{S,S'}\prod_{i\in S}\tau_i^{\mathrm{in}}\prod_{i\in S'}y_i$ can be computed by multiplying together one subset from each group, namely the restriction of $S$ to the group. Instead of $2^n$ registers for all subsets, we use only $k\cdot2^{\lceil n/k\rceil}$ registers corresponding to subsets in the $k$ groups, and we can compute all the corresponding monomials into these registers in time $2n$ using the first half of Algorithm1. Then since we are only multiplying $k$ monomials together, we can compute $q_f$ using Algorithm2 in time $2^k\cdot2\cdot2n$, since each call to ToggleSomeInputs is replaced with our $2n$ time execution of Algorithm1.
As a slight last speedup, we use a Gray code to order our subsets in Algorithm2, replacing two executions of ToggleSomeInputs with a subroutine toggling a single group on or off and only spending $2\lceil n/k\rceil$ time to do so. This allows us to compute $q_f$ in $2^k\cdot2\lceil n/k\rceil$ time rather than $2^k\cdot4n$ time.
\begin{proof}[Proof of Theorem11]
For $j\in[k]$ let $b_j=\lceil nj/k\rceil$, and divide the range $[n]$ into $k$ groups: $G_1=\{1,\ldots,b_1\}$, $G_2=\{b_1+1,\ldots,b_2\}$, $\ldots$, $G_k=\{b_{k-1}+1,\ldots,b_n=n\}$. For each group $G_j$, we have $2^{|G_j|}$ registers $R_{j,S}$ indexed by subsets $S\subseteq G_j$. As in all previous algorithms we also use $n$ registers $R_1^{\mathrm{in}},\ldots,R_n^{\mathrm{in}}$, corresponding to the $n$ input bits, for a total of $n+\sum_{j=1}^k2^{|G_j|}$ registers plus the output register $R_{\mathrm{out}}$.
We'll begin by rewriting the polynomial $q_f$ in a new form. Recall from Section3.1 that $q_f(\vec\tau^{\mathrm{in}},\vec y)=f(x)$ so long as $\vec y=\vec x+\vec\tau^{\mathrm{in}}$, and
\[q_f(\vec\tau^{\mathrm{in}},\vec y)=\sum_{S,S'\subseteq[n]}c_{S,S'}\left(\prod_{i\in S}\tau_i^{\mathrm{in}}\right)\left(\prod_{i\in S'}y_i\right).\]
Now, for every $S'\subseteq[n]$, define $S'_j:=S'\cap G_j$. For each $j\in[k]$ and $S\subseteq G_j$, let $z_{j,S}=\tau_{j,S}+\prod_{i\in S}y_i$, where $\tau_{j,S}$ is the initial value of register $R_{j,S}$. Now for every monomial in $q_f$, we split the term $\prod_{i\in S'}y_i$ in the monomial into $k$ different products $\prod_{i\in S'_j}y_i$, each of which we can replace with $z_{j,S'_j}-\tau_{j,S'_j}$. This gives us a new polynomial
\[r_f(\vec\tau^{\mathrm{in}},\vec\tau,\vec z)=\sum_{S,S'\subseteq[n]}c_{S,S'}\left(\prod_{i\in S}\tau_i^{\mathrm{in}}\right)\left(\prod_{j=1}^k(z_{j,S'}-\tau_{j,S'})\right).\]
As we did with $q_f$, for $S,S'\subseteq[n]$ and $T\subseteq[k]$, let $d_{S,S',T}$ be the coefficient of $(\prod_{i\in S}\tau_i^{\mathrm{in}})\cdot(\prod_{j\in T}z_{j,S'_j})\cdot(\prod_{j\in[k]\setminus T}\tau_{j,S'_j})$ in $r_f$, so that
\begin{multline*}
r_f(\vec\tau^{\mathrm{in}},\vec\tau,\vec z)=\sum_{S\subseteq[n]}\sum_{S'\subseteq[n]}\sum_{T\subseteq[k]}d_{S,S',T}\\
\cdot\left(\prod_{i\in S}\tau_i^{\mathrm{in}}\right)\left(\prod_{j\in T}z_{j,S'_j}\right)\left(\prod_{j\in[k]\setminus T}\tau_{j,S'_j}\right)
\end{multline*}
which is equivalent to $f(x_1\ldots f_n)$ as long as $z_{j,S'_j}=\tau_{j,S'_j}+\prod_{i\in S'_j}(x_i+\tau_i^{\mathrm{in}}+1)$.
Following ToggleSomeInputs($S'$), we define new versions of ToggleInput and ToggleMonomials from Section3.1 which focus on some groups and not others. In fact we will only focus on a single group $G_j$ rather than a subset of the groups, as we will order our subsets $S'$ in such a way that we will only ever need to toggle one group at a time:
\paragraph{ToggleInputForGroup($j$).}
\begin{enumerate}[label=\arabic*:]
\item \textbf{procedure} ToggleInputForGroup($j$)
\item \textbf{for} $i\in G_j$ \textbf{do}
\item $R_i^{\mathrm{in}}\leftarrow R_i^{\mathrm{in}}+x_i$
\item \textbf{end for}
\item \textbf{end procedure}
\end{enumerate}
\paragraph{ToggleMonomialsForGroup($j$).}
\begin{enumerate}[label=\arabic*:]
\item \textbf{procedure} ToggleMonomialsForGroup($j$)
\item ToggleInputForGroup($j$)
\item \textbf{for} $S\subseteq G_j$ \textbf{do}
\item $R_{j,S}\leftarrow R_{j,S}+\prod_{i\in S}R_i^{\mathrm{in}}$
\item \textbf{end for}
\item ToggleInputForGroup($j$)
\item \textbf{end procedure}
\end{enumerate}
We are now ready to assemble Algorithm4, which completes the proof of Theorem11. To incorporate our improvement using Gray codes[9], let $T_0=\emptyset,\ldots,T_{2^k-1}$ be an ordering of all subsets of $[k]$ such that each consecutive pair of sets $T_{\ell},T_{\ell+1\bmod2^k}$ differs by exactly one element $e_{\ell}\in[k]$; thus we will only need to toggle the group $G_{e_{\ell}}$ as claimed:
\paragraph{Algorithm3: Transparently compute $f$ in time $2^k\cdot2\lceil n/k\rceil$ with $n+k\cdot2^{\lceil n/k\rceil}$ registers.}
\begin{enumerate}[label=\arabic*:]
\item \textbf{for} $\ell=0,\ldots,2^k-1$ \textbf{do}
\item \[\begin{aligned}R_{\mathrm{out}}\leftarrow R_{\mathrm{out}}+\sum_{S\subseteq[n]}\sum_{S'\subseteq[n]}&d_{S,S',T_{\ell}}\\
&\mathord{}(\prod_{i\in S}R_i^{\mathrm{in}})(\prod_{j=1}^kR_{j,S'_j})\end{aligned}\]
\item ToggleMonomialsForGroup($e_{\ell}$)
\item \textbf{end for}
\end{enumerate}
Each time Line2 is reached, we have $R_{j,S}=\tau_{j,S}+\prod_{i\in S}y_j$ for $j\in T_{\ell}$, and $R_{j,S}=\tau_{j,S}$ for $j\in[k]\setminus T_{\ell}$. We also have $R_i^{\mathrm{in}}=\tau_i^{\mathrm{in}}$ for each $i\in[n]$. So the effect of the line is to add
\[\sum_{S\subseteq[n]}\sum_{S'\subseteq[n]}d_{S,S',T_{\ell}}\left(\prod_{i\in S}\tau_i^{\mathrm{in}}\right)\left(\prod_{j\in T_{\ell}}z_{j,S'_j}\right)\left(\prod_{j\in[k]\setminus T_{\ell}}\tau_{j,S'_j}\right)\]
to $R_{\mathrm{out}}$. Summing this expression over all possible subsets $T_{\ell}\subseteq[k]$ gives $R_{\mathrm{out}}=\tau^{\mathrm{out}}+r_f(\cdots)=\tau^{\mathrm{out}}+f(x_1,\ldots,x_n)$, and so our algorithm transparently computes $f$.
\end{proof}
\paragraph{Note12.} It is not difficult to save $k$ registers by removing $R_{1,\emptyset},\ldots,R_{k,\emptyset}$, as we simply add each value $d_{S,\emptyset,T}(\prod_{i\in S}R_i^{\mathrm{in}})$ to our polynomial without concerning ourselves with any $x_i$ (and by extension any $y_i$ or $z_i$) variables.
